% ===========================================================================
\section{Shared Nitrogen Accounting and Air-Pollutant Methods}
\label{sec:methods-shared}
% ===========================================================================

The greenhouse-gas and air-pollutant calculations use a common nitrogen ledger. This ledger
tracks how excreted and applied nitrogen moves through manure management, grazing,
fertiliser application, volatilisation and soil processes.

The two inventories therefore represent different emissions from the same underlying
nitrogen flows. Aligning their system boundaries is necessary to prevent the same nitrogen
from being removed twice, retained twice, or assigned inconsistently between \NHthree{} and
\NtwoO{}.

\subsection{Alignment of the nitrogen boundaries}
\label{sec:methods-ledger}

The reference TAN formulation follows species- and system-specific ammonia loss through housing or yard exposure, storage and field application. Each stage acts on surviving total ammoniacal nitrogen. Management NH$_3$ and NO$_x$, leaching, N$_2$ and N$_2$O losses are removed once before the surviving organic pool reaches soil. Grazing and field volatilisation remain separate pathways. The same realised pools drive indirect N$_2$O and the air species.

A separately identified IPCC alternative uses the applicable management and field volatilisation fractions. It regenerates both product families rather than combining an IPCC N$_2$O account with a different EMEP ammonia chain. Its management fractions are not additional removal from the TAN reference. Material-flow closure is tested in each configuration and uncertainty draw.

\subsection{Synthetic-fertiliser ammonia}
\label{sec:methods-air}

Synthetic-fertiliser \NHthree{} is calculated from Pakistan's reported fertiliser-product
mix rather than from one generic fertiliser composition.

EMEP/EEA factors vary substantially among products because urea and ammonium-based
fertilisers have greater volatilisation potential than nitrate-dominant products,
particularly on high-pH soils. Pakistan's alkaline cropland and urea-heavy fertiliser mix
make product weighting important.

Let $w_j$ be product $j$'s share of applied nitrogen and $EF^{\NHthree{}\text{-}N}_j$ its
ammonia-nitrogen emission factor \citep{emep2023soils}. Taking the NFDC offtake by product
\citep{nfdc2024} over one fiscal year, the product-weighted factor is
\begin{equation}
f_{\NHthree{},SF} \;=\; \sum_j w_j\, EF^{\NHthree{}\text{-}N}_j \;=\; 0.1608 .
\end{equation}
Urea supplies approximately \SI{81}{\percent} of the nitrogen in the product mix and
therefore dominates the weighted result (Table~\ref{tab:fertmix}). The remaining products
change the factor by less than 0.01.

The value 0.1608 is already expressed as kilograms of \NHthree{}-N per kilogram of nitrogen
applied. It is converted to \NHthree{} mass by the ratio $17/14$ and used directly in the
ammonia calculation. It is not multiplied by an additional \NHthree{} share.

Soil \NOx{} is calculated separately, using $f_{NO,SF}=0.007$ kg NO-N per kilogram of
nitrogen applied. For the IPCC indirect-\NtwoO{} pathway, which Eq.~11.9 defines on the
combined volatilised pool, the two are summed:
\begin{equation}
\mathrm{Frac_{GASF}} \;=\; f_{\NHthree{},SF} + f_{NO,SF} \;=\; 0.1678 ,
\end{equation}
against the IPCC ``urea-dominant'' default of 0.15. Keeping these three quantities separate
preserves the nitrogen balance and avoids applying an additional ammonia fraction to a
factor already expressed as \NHthree{}-N. Because $f_{\NHthree{},SF}$ also drives the
greenhouse-gas layer's synthetic indirect \NtwoO{}, it is resolved in one place and both
layers necessarily use it.

\begin{table*}[!htbp]
\centering
\footnotesize
\begin{tabular}{@{}lrrrrr@{}}
\toprule
\textbf{Product} & \textbf{t N} & \textbf{N content} & \textbf{Share of N}
& \textbf{$EF^{\NHthree{}\text{-}N}$} & \textbf{Contribution} \\
\midrule
Urea             & 3\,026\,265 & 0.46 & 0.8078 & 0.170 & 0.1373 \\
DAP              &   290\,777 & 0.18 & 0.0776 & 0.154 & 0.0120 \\
CAN              &   206\,058 & 0.26 & 0.0550 & 0.043 & 0.0024 \\
Other NP         &   202\,590 & 0.20 & 0.0541 & 0.154 & 0.0083 \\
Ammonium sulphate &   14\,595 & 0.21 & 0.0039 & 0.154 & 0.0006 \\
NPK              &     6\,162 & 0.15 & 0.0016 & 0.154 & 0.0003 \\
\midrule
\textbf{Total}   & \textbf{3\,746\,448} & & \textbf{1.0000} & & \textbf{0.1608} \\
\bottomrule
\end{tabular}
\caption{Derivation of the product-weighted synthetic-fertiliser ammonia factor. The table is
generated from \texttt{outputs/agriculture\_fertiliser\_composition\_2024.csv}, whose window is
exactly the twelve months July~2023 to June~2024; the annexe supplies the normalised composition
only, and the nitrogen magnitude remains the audited national offtake. NFDC's national and
provincial tables do not sum to each other over the same twelve months: the difference is
\SI{15918}{\tonne} of product, \SI{0.154917}{\percent} of the national total and mostly
urea. It is published by month and product in
\texttt{outputs/agriculture\_nfdc\_province\_residual.csv} and assigned to
Other/unallocated rather than distributed across the provinces, which would make every
residual zero and the shares look exact. Product tonnages are
NFDC offtake over one fiscal year \citep{nfdc2024}; nitrogen contents and emission factors are
EMEP/EEA values for high-pH soils, appropriate to Pakistan's alkaline cropland
\citep{emep2023soils}. \emph{All emission factors in this table are kg \NHthree{}-N per kg N
applied.} Urea carries \SI{81}{\percent} of the nitrogen and therefore dominates the weighting; the
audited national nitrogen total (\SI{3.775}{}~Mt~N) rather than this window's sum is used as the
activity datum, with the mix supplying only the composition.}
\label{tab:fertmix}
\end{table*}

\subsection{Livestock ammonia}

Livestock \NHthree{} is calculated using the EMEP/EEA total-ammoniacal-nitrogen (TAN)
mass-flow method \citep{emep2023manure,sommer2019}. The IPCC post-application fraction
$\mathrm{Frac_{GASM}}$ is not used as a livestock ammonia emission factor, because it was
developed for indirect-\NtwoO{} accounting rather than stage-resolved \NHthree{} estimation.

Excreta deposited during grazing lose a species-specific proportion of TAN directly. Manure
entering a managed system can lose ammonia during housing or yard exposure, storage and
field application. Each stage acts on the TAN remaining after the preceding stage.

The calculation distinguishes among manure systems. Solid storage passes through housing,
storage and field-application stages. Daily-spread manure reaches the field without a
storage stage. Dry lots are represented as uncovered yards.

Factors also differ among livestock categories, and are taken from Table~3-9 of the
Guidebook. Buffalo manure is assigned a TAN fraction of \num{0.50}, compared with
\num{0.60} for cattle, and a storage factor of \num{0.17} against \num{0.32}. Grazing
factors are applied to TAN rather than to total excreted nitrogen (\num{0.14} for cattle and
buffalo, \num{0.09} for small ruminants).

EMEP/EEA does not provide a separate camel category. Camels therefore use the
small-ruminant factors as the closest available published analogue, with a wider uncertainty
assigned to reflect the approximation.

\subsection{Other agricultural air pollutants}

Soil \NOx{} is estimated as approximately \SI{0.7}{\percent} of applied synthetic, organic
and crop-residue nitrogen, following \citet{stehfest2006}. The value is initially calculated
as NO-N and converted to reported \NOx{} mass using the conventional NO$_2$-equivalent
conversion.

Pollutants from crop-residue burning are calculated from the same dry-matter-burned estimate
used in the greenhouse-gas inventory. The burned mass is multiplied by the combustion factor
and by pollutant-specific emission factors for CO, PM$_{2.5}$, NMVOC, \NOx{}, \NHthree{},
black carbon and SO$_2$ (Table~\ref{tab:apef}). The PM$_{2.5}$ factor is also used by the
related satellite-nowcasting layer so that the two products report a consistent crop-burning
PM$_{2.5}$ quantity.

All air factors retain native units and source applicability. Local product-rule widths are diagnostics; the integrated conditional ensemble replays shared herd, crop, nutrient and bounded fate-factor realisations across NH$_3$, NO$_x$ and N$_2$O. Structural management shares and factor alternatives are separate cases, not confidence limits.

\begin{table}[!htbp]
\centering
\footnotesize
\setlength{\tabcolsep}{3pt}
\begin{tabular}{@{}p{0.26\linewidth}p{0.40\linewidth}rr@{}}
\toprule
\textbf{Pollutant / stage} & \textbf{Basis} & \textbf{Factor} & \textbf{$\pm$\%} \\
\midrule
\NHthree{}, synthetic fert.\ & $f_{\NHthree{},SF}$, kg \NHthree{}-N per kg N & 0.1608$^{b}$ & 30 \\
\NHthree{}, grazing & of TAN; cattle/buffalo \textbar{} small rum. & 0.14 \textbar{} 0.09 & 50 \\
\NHthree{}, housing & of TAN; cattle \textbar{} buffalo \textbar{} small rum. & 0.08 \textbar{} 0.20 \textbar{} 0.22 & 40 \\
\NHthree{}, yard (dry lot) & of TAN; dairy \textbar{} other cattle & 0.30 \textbar{} 0.53 & 50 \\
\NHthree{}, storage & of TAN surviving; cattle \textbar{} buffalo & 0.32 \textbar{} 0.17 & 40 \\
\NHthree{}, field application & of TAN surviving; cattle \textbar{} buffalo & 0.68 \textbar{} 0.55 & 40 \\
TAN fraction & of N excreted; cattle \textbar{} buffalo \textbar{} poultry & 0.60 \textbar{} 0.50 \textbar{} 0.70 & 20 \\
Soil \NOx{} & $f_{NO,SF}$, kg NO-N per kg N applied$^{a}$ & 0.007 & 100 \\
\multicolumn{4}{@{}l}{\emph{Crop-residue burning (g/kg dry matter; on $M_{burnt}\!\cdot\!C_f$)}} \\
\quad CO       & IPCC 2006 T2.5      & 92 & 40 \\
\quad PM$_{2.5}$ & \citet{akagi2011}  & 7.6 & 50 \\
\quad NMVOC    & \citet{akagi2011} / EMEP & 7.0 & 60 \\
\quad \NOx{}   & IPCC 2006 T2.5      & 2.5 & 50 \\
\quad \NHthree{} & \citet{akagi2011} & 1.3 & 60 \\
\quad BC       & \citet{akagi2011}   & 0.5 & 60 \\
\quad SO$_2$   & \citet{emep2019}    & 0.40 & 70 \\
\bottomrule
\end{tabular}
\caption{Air-pollutant factors, calculation basis, and 95\% relative uncertainty. Livestock
\NHthree{} follows the species-resolved EMEP/EEA TAN mass-flow method. Crop-burning factors are
applied to the same dry matter used in the greenhouse-gas calculation. The synthetic-fertiliser
\NHthree{} factor is already expressed as \NHthree{}-N and receives no additional
ammonia-share adjustment. Soil \NOx{} is calculated as NO-N and reported as
NO$_2$-equivalent mass.}
\label{tab:apef}
\end{table}
